\chapter[Presentacion]{Presentación}
\label{cp:introduction}

\insertminitoc
\parindent0pt

\section{Planteamiento del Problema}
El uso de lo que conocemos como dispositivos móviles por parte de los niños y adolescentes ha aumentado de forma desmedida en las últimas décadas, convirtiéndose en un de los elemento habitual de su vida cotidiana para el entretenimiento, la comunicación y el aprendizaje. De acuerdo con el informe \cite{noauthor_full_2025} una proporción significativa de niños accede a internet desde edades tempranas y alrededor del 70\,\% de los niños de aproximadamente 10 años ya cuenta y dispone de un teléfono inteligente propio en países miembros de la Organización para la Cooperación y el Desarrollo Económicos.\\

Esta exposición temprana y prolongada al entorno digital implica diversos riesgos para la población infantil, tales como el uso excesivo de dispositivos, la exposición a contenidos no adecuados para su edad y la participación en interacciones digitales que pueden resultar perjudiciales. La literatura científica señala que el uso intensivo de tecnologías digitales puede generar efectos negativos a nivel emocional y psicológico en los niños, incluyendo ansiedad, alteraciones del estado de ánimo y dificultades en la regulación emocional, especialmente cuando no existe una supervisión adulta adecuada \cite{limone_psychological_2021}.\\

Aunque padres y tutores reconocen la importancia de supervisar el uso de la tecnología en la infancia, esta labor resulta compleja debido a limitaciones de tiempo, desconocimiento del contenido consumido y a la naturaleza personal del uso de los dispositivos móviles. Como consecuencia, numerosas situaciones de riesgo digital no son detectadas de manera oportuna, sino hasta que sus efectos ya se manifiestan en el comportamiento, el rendimiento académico o el bienestar emocional del menor \cite{noauthor_full_2025}.\\

Contemplando este contexto, se evidencia una problemática relevante relacionada con la creciente exposición de niños y adolescentes a entornos digitales potencialmente riesgosos, sin mecanismos suficientes que permitan comprender y dimensionar las situaciones a las que están expuestos durante el uso cotidiano de dispositivos móviles, lo que representa un desafío significativo para la protección integral de la infancia en la sociedad digital contemporánea.\\

\subsection{Preguntas de Investigación}

\begin{itemize}
   \item ¿Cuáles son los principales riesgos digitales a los que están expuestos los niños durante el uso cotidiano de dispositivos móviles?
    
    \item ¿Qué tipos de contenidos e interacciones digitales representan una mayor amenaza para el bienestar emocional y psicológico de los niños?
    
    \item ¿De qué manera el tiempo y la frecuencia de uso de dispositivos móviles influyen en la aparición de situaciones de riesgo digital en la población infantil?
    
    \item ¿Qué dificultades enfrentan los padres o tutores para identificar de manera oportuna los riesgos digitales presentes en el uso de dispositivos móviles por parte de los niños?
    
    \item ¿Cómo se manifiestan las situaciones de riesgo digital en el comportamiento y bienestar de los niños durante su interacción con el entorno digital?
    
    \item ¿Qué limitaciones existen actualmente para la comprensión y el registro de los riesgos digitales asociados al uso infantil de dispositivos móviles?
\end{itemize}

\subsection{Justificación del Estudio}
El crecimiento acelerado del uso de dispositivos móviles en la infancia constituye una problemática relevante a nivel global, debido a su impacto potencial en el desarrollo emocional, psicológico y social de los niños. En la actualidad, los dispositivos móviles forman parte integral de la vida cotidiana de los menores, facilitando el acceso a información, comunicación y entretenimiento; sin embargo, esta exposición temprana y prolongada al entorno digital también incrementa la probabilidad de que los niños se enfrenten a contenidos inapropiados y a interacciones digitales de riesgo. La falta de comprensión profunda sobre estas situaciones limita la capacidad de los padres, educadores y responsables institucionales para actuar de manera preventiva y oportuna, lo que justifica la necesidad de estudiar esta problemática desde una perspectiva académica y social.

A nivel local, en el municipio de Comayagua, el acceso a teléfonos inteligentes y a internet móvil ha aumentado considerablemente en los últimos años, incluso en hogares donde la supervisión adulta es limitada debido a factores laborales, económicos o educativos. En muchos casos, los dispositivos móviles se utilizan como una herramienta de entretenimiento o apoyo educativo para los niños, sin un acompañamiento constante que permita identificar los riesgos asociados al uso inadecuado del entorno digital. Esta realidad local refleja una problemática que, aunque forma parte de una tendencia global, adquiere características particulares en contextos donde la alfabetización digital y los mecanismos de supervisión parental son insuficientes.

Desde el ámbito académico, esta investigación resulta relevante porque contribuye a la comprensión de los riesgos digitales presentes en el uso infantil de dispositivos móviles, abordando un tema de creciente interés en las áreas de tecnología, educación y bienestar infantil. El estudio permite generar conocimiento que puede servir como base para futuras investigaciones, políticas públicas o estrategias educativas orientadas a la protección de la infancia en entornos digitales, tanto a nivel local como nacional.

Asimismo, la investigación se justifica desde una perspectiva social, ya que busca visibilizar una problemática que afecta directamente a niños, padres y tutores, promoviendo una mayor conciencia sobre la importancia de la supervisión y el acompañamiento en el uso de tecnologías digitales. Comprender esta realidad es un paso fundamental para fortalecer la protección de los derechos de la niñez y fomentar un uso más responsable y seguro de los dispositivos móviles en la sociedad contemporánea.
\begin{itemize}
    \item \textbf{Social}, al contribuir a la protección del bienestar emocional y psicológico de los niños mediante una mejor comprensión de los riesgos digitales a los que están expuestos durante el uso de dispositivos móviles.
    
    \item \textbf{Familiar y educativa}, al apoyar a padres, tutores y educadores en la identificación de situaciones digitales potencialmente riesgosas, fortaleciendo la supervisión y el acompañamiento en el entorno digital infantil.
    
    \item \textbf{Académica}, al generar conocimiento relevante sobre los riesgos digitales en la infancia, sirviendo como base para futuras investigaciones, estrategias educativas y políticas orientadas a la protección de la niñez en contextos digitales.
\end{itemize}

\subsection{Viabilidad del Estudio}

La viabilidad del presente proyecto se fundamenta en la convergencia de la necesidad social de protección infantil y la disponibilidad actual de tecnologías de inteligencia artificial de alto rendimiento. A continuación, se detalla el análisis basado en los recursos, alcances e implicaciones.

\begin{enumerate} 
    \item \textbf{Disponibilidad de los recursos}
    
    \begin{itemize} 
        \item \textbf{Humanos:} La investigación y el desarrollo del prototipo funcional serán ejecutados por el tesista, contando con la dirección especializada de asesores en las áreas de ingeniería de software y seguridad informática. Se integran además conocimientos transversales de psicología infantil para asegurar que la herramienta cumpla con los estándares de mediación parental habilitadora \cite{livingstone_maximizing_2017, mcnally_co-designing_2018}.
                
        \item \textbf{Datos:} Para el entrenamiento y validación de la \gls{ia}, se emplearán bases de datos públicas y repositorios de investigación que contienen patrones de lenguaje y etiquetas de contenido sensible. Ante la naturaleza privada de las interacciones, se utilizarán técnicas de generación de datos sintéticos y marcos de trabajo existentes para la detección de riesgos como el \textit{grooming} y el \textit{cyberbullying} \cite{razi_sliding_2023, zhu_cyberbullying_2021}.
    \end{itemize}

    \item \textbf{Alcances del estudio}
    
    La investigación se delimita al diseño, desarrollo y pruebas de concepto de una aplicación móvil para el sistema operativo Android. El sistema se centrará en la detección de tres categorías críticas de riesgo: conversaciones con contenido inapropiado, interacciones con desconocidos y visualización de imágenes sensibles. El estudio no busca reemplazar la comunicación familiar, sino proporcionar una "bitácora de evidencia" que sirva de insumo para la mediación parental activa, demostrando que la inteligencia artificial puede actuar como un observador no intrusivo que protege la integridad del menor en entornos de mensajería y redes sociales \cite{stoilova_parental_2024, rojas_technical_2025}.

    \item \textbf{Implicaciones del estudio} 
    
    \begin{itemize}
        \item \textbf{Académicas:} La tesis contribuye al campo de la seguridad informática y la interacción humano-computadora en Honduras, proporcionando un precedente sobre el uso de visión artificial aplicada a la protección de grupos vulnerables en entornos digitales.
        
        \item \textbf{Tecnológicas:} Proporciona una arquitectura metodológica para la supervisión proactiva que supera las limitaciones de los filtros de DNS y listas negras tradicionales, los cuales han demostrado ser ineficaces y fáciles de evadir en el contexto móvil actual.
        
        \item \textbf{Sociales:} Ofrece una herramienta que empodera a los padres de familia, reduciendo la brecha digital y mitigando los efectos negativos de la exposición no supervisada, tales como la ansiedad y el distrés psicológico derivados del ciberacoso.
    \end{itemize}
    
    \item \textbf{Consecuencias del estudio}
    
    Los resultados que se van a presentar en esta investigación permitirán validar la hipótesis de que es posible supervisar de forma automatizada y privada el contenido de un menor sin violar su autonomía, siempre que exista un marco de confianza mediado por el acceso parental mediante PIN secreto. La metodología y el algoritmo de detección resultantes tienen el potencial de ser escalados a soluciones comerciales o integrados en políticas de bienestar digital, sirviendo como modelo para otras aplicaciones que busquen proteger a los niños y adolescentes en la era post-pandemia \cite{limone_psychological_2021, noauthor_full_2025}.
\end{enumerate}

\section{Objetivos de la Investigación}
\subsection{Objetivo General}
Desarrollar un Sistema de Control Parental Basado en Aprendizaje Automático orientado a la Protección Infantil en la Era Digital
\subsection{Objetivos Especificos}
\begin{enumerate}
    \item Analizar patrones de uso del dispositivo móvil infantil para identificar contextos digitales asociados a riesgos, tales como conversaciones sensibles, interacción con desconocidos y exposición a contenido inapropiado.
    \item Definir cuales son los criterios de supervisión parental más relevantes para la protección infantil en el entorno digital, considerando aspectos como la edad del niño, el tipo de contenido consumido y las interacciones digitales realizadas.
    \item Implementar un prototipo funcional de un sistema de apoyo a la supervisión parental basado en aprendizaje automático y eficiente.
    \item Evaluar la precisión y efectividad del modelo inteligente en la detección de eventos de riesgo, así como su impacto en el rendimiento del dispositivo y la experiencia del usuario.
\end{enumerate}